\documentclass[12pt, a4paper]{article}

% --- Pacotes Básicos e Normas ABNT (Portáveis e Livres de Erros) ---
\usepackage[utf8]{inputenc}
\usepackage[T1]{fontenc}
\usepackage[margin=2.5cm]{geometry}
\usepackage{amsmath, amsfonts, amssymb, amsthm}
\usepackage{graphicx}
\usepackage{booktabs}
\usepackage{cite}
\usepackage{url}
\usepackage{listings}
\usepackage{xcolor}
\usepackage{tikz}
\usetikzlibrary{graphs, positioning, shapes.geometric}
\usepackage[brazilian]{babel}
\usepackage{ragged2e}


\renewcommand{\abstractname}{Resumo}

% --- Configuração de Código Python (Listings) ---
\definecolor{codegreen}{rgb}{0,0.6,0}
\definecolor{codegray}{rgb}{0.5,0.5,0.5}
\definecolor{codepurple}{rgb}{0.58,0,0.82}
\definecolor{backcolour}{rgb}{0.96,0.96,0.96}

\lstdefinestyle{pythonstyle}{
    backgroundcolor=\color{backcolour},   
    commentstyle=\color{codegreen},
    keywordstyle=\color{magenta},
    numberstyle=\tiny\color{codegray},
    stringstyle=\color{codepurple},
    basicstyle=\ttfamily\footnotesize,
    breakatwhitespace=false,         
    breaklines=true,                 
    captionpos=b,                    
    keepspaces=true,                 
    numbers=left,                    
    numbersep=5pt,                  
    showspaces=false,                
    showstringspaces=false,
    showtabs=false,                  
    tabsize=4
}
\lstset{style=pythonstyle}

% --- Elementos Pré-textuais ---
\title{\Large\bfseries \justifying Otimização Combinatória e Teoria dos Grafos Aplicadas ao Sudoku: Um Estudo da Coloração Própria e Algoritmos de Backtracking Explicitados em Python no Ensino de Matemática}

\author{\textbf{Joao Victor Silva de Sousa}\\[0.5em]
\small Curso de Licenciatura em Matemática\\
\small Disciplina: Metodologia da Pesquisa Educacional III\\
\small Instituto Federal do Amazonas}

\date{\today}

\begin{document}

\maketitle

\begin{abstract}
Este artigo aborda a resolução e a geração automatizada de enigmas do jogo Sudoku sob a ótica da Matemática Aplicada, Otimização Combinatória e Teoria dos Grafos. O tabuleiro $9 \times 9$ é formalizado rigorosamente como um grafo simples não orientado $G = (V, E)$ com $|V| = 81$ vértices e $|E| = 810$ arestas, onde a busca por uma solução equivale a encontrar uma $9$-coloração própria para o grafo. Desenvolve-se em linguagem Python uma implementação do algoritmo de \textit{backtracking} que opera diretamente sobre o conjunto de adjacências $N(v)$ do grafo explícito, superando a matriz convencional $9 \times 9$, e um algoritmo de remoção progressiva com validação de unicidade de solução. Discute-se a relevância pedagógica desta modelagem no âmbito da formação inicial de professores de Matemática, articulando os pilares do Pensamento Computacional e as competências estipuladas pela Base Nacional Comum Curricular (BNCC), sob a normatização técnica da ABNT NBR 14724.
\\[0.5em]
\textbf{Palavras-chave:} Otimização Combinatória; Teoria dos Grafos; Coloração Própria; Backtracking em Grafos; Python; BNCC; ABNT.
\end{abstract}

\vspace{1em}
\noindent\rule{\linewidth}{0.4pt}
\vspace{1.5em}

\section{Introdução}
No âmbito da Matemática Aplicada e da Otimização Combinatória, a formalização de jogos e quebra-cabeças lógicos constitui um excelente ambiente para o desenvolvimento do raciocínio formal e algorítmico. O Sudoku, um enigma combinatório popular, é frequentemente resolvido na prática por métodos empíricos e heurísticas de eliminação visual. Contudo, sob uma perspectiva matemática rigorosa, a resolução do Sudoku pertence à classe dos problemas de Satisfabilidade de Restrições (CSP) e pode ser formulada como o problema da $k$-coloração própria de um grafo.

Na formação inicial de docentes de Matemática, a integração entre conteúdos de Matemática Discreta e ferramentas computacionais atende diretamente às diretrizes da Base Nacional Comum Curricular (BNCC), especialmente no que tange ao desenvolvimento do \textbf{Pensamento Computacional} e da \textbf{Modelagem Matemática} (habilidades EM13MAT315 e EM13MAT405). 

Este trabalho apresenta a modelagem matemática completa do grafo do Sudoku, a implementação explícita do algoritmo de \textit{backtracking} fundamentado no conjunto de vizinhanças $N(v)$ em linguagem Python, e o algoritmo de escavação com garantia de unicidade de solução para geração de tabuleiros, estruturado segundo os ditames da ABNT NBR 14724.

\section{Modelagem Matemática: O Grafo do Sudoku}

\subsection{Definição do Grafo $G = (V, E)$}
Definimos o tabuleiro de Sudoku $9 \times 9$ como um grafo simples, não orientado e não ponderado $G = (V, E)$, definido pelas seguintes estruturas:
\begin{itemize}
    \item \textbf{Conjunto de Vértices ($V$):} Cada uma das 81 células corresponde a um vértice $v_{i,j} \in V$, com $i, j \in \{1, 2, \dots, 9\}$. Assim, $|V| = 81$. Indexando linearmente os vértices para implementação computacional, temos $k = 9(i-1) + (j-1)$, com $k \in \{0, 1, \dots, 80\}$.
    \item \textbf{Conjunto de Arestas ($E$):} Uma aresta não orientada $(u, w) \in E$ existe entre dois vértices distintos $u = v_{i,j}$ e $w = v_{k,l}$ se e somente se eles compartilharem ao menos uma relação de conflito.
\end{itemize}

As relações de conflito são dadas pelas proposições lógicas:
\begin{enumerate}
    \item \textbf{Conflito de Linha:} $i = k$;
    \item \textbf{Conflito de Coluna:} $j = l$;
    \item \textbf{Conflito de Subgrade $3 \times 3$:} $\lfloor \frac{i-1}{3} \rfloor = \lfloor \frac{k-1}{3} \rfloor \land \lfloor \frac{j-1}{3} \rfloor = \lfloor \frac{l-1}{3} \rfloor$.
\end{enumerate}

Portanto, o conjunto de arestas é formalmente expresso por:
\begin{equation}
	\begin{aligned}
		E = \Bigl\{ (v_{i,j}, v_{k,l}) \in V \times V \Bigm| &\ (i,j) \neq (k,l) \ \land \ \Bigl( i=k \ \lor \ j=l \ \lor \\ 
		&\ \Bigl( \Big\lfloor\frac{i-1}{3}\Big\rfloor = \Big\lfloor\frac{k-1}{3}\Big\rfloor \ \land \ 
		\Big\lfloor\frac{j-1}{3}\Big\rfloor = \Big\lfloor\frac{l-1}{3}\Big\rfloor \Bigr) \Bigr) \Bigr\}
	\end{aligned}
\end{equation}

\subsection{Propriedades Topológicas e Grau dos Vértices}
Seja $v \in V$ um vértice genérico. Sua vizinhança aberta $N(v)$ é o conjunto de vértices adjacentes a $v$:
\begin{equation}
    N(v) = \{u \in V \mid (u,v) \in E\}
\end{equation}

Para calcular o grau $d(v) = |N(v)|$, contabilizam-se as adjacências sem sobreposição:
\begin{itemize}
    \item $8$ vértices na mesma linha;
    \item $8$ vértices na mesma coluna;
    \item $4$ vértices adicionais no mesmo bloco $3 \times 3$ que não pertencem à mesma linha nem à mesma coluna.
\end{itemize}

Assim, $G$ é um grafo $20$-regular, ou seja, $d(v) = 20, \forall v \in V$. Pelo Lema dos Apertos de Mão (Handshaking Lemma):
\begin{equation}
    |E| = \frac{1}{2} \sum_{v \in V} d(v) = \frac{81 \times 20}{2} = 810 \text{ arestas.}
\end{equation}

\subsection{Formulação como $9$-Coloração Própria de Grafos}
Uma coloração de $G$ utilizando $k = 9$ cores (correspondentes aos dígitos $\{1, 2, \dots, 9\}$) é uma função $c: V \to \{1, 2, \dots, 9\}$. A coloração é dita \textbf{própria} se:
\begin{equation}
    \forall (u,w) \in E \implies c(u) \neq c(w)
\end{equation}

Dado um estado inicial com pistas pré-definidas $V_{fixo} \subset V$ e uma atribuição inicial $c_0: V_{fixo} \to \{1, \dots, 9\}$, a resolução do Sudoku consiste em estender $c_0$ para uma coloração própria total $c: V \to \{1, \dots, 9\}$.

\section{Implementação Explícita do Algoritmo de Backtracking em Grafo}

A maioria das implementações computacionais tradicionais valida restrições varrendo linhas e colunas através de matrizes bidimensionais. Nesta abordagem de Matemática Aplicada, o algoritmo constrói \textbf{explicitamente} a estrutura topológica do grafo $G=(V,E)$ como uma lista de adjacência $N(v)$ e executa o \textit{backtracking} consultando diretamente a vizinhança topológica.

\subsection{Cálculo do Conjunto de Cores Viáveis por Operação de Conjuntos}
Dado o estado corrente de coloração $c$, o conjunto de cores indisponíveis para o vértice $v$ é a imagem do seu conjunto vizinho $N(v)$:
\begin{equation}
    C_{indisponivel}(v) = \{ c(u) \in \{1, \dots, 9\} \mid u \in N(v) \land c(u) \neq 0 \}
\end{equation}

O conjunto de cores viáveis $C_{viavel}(v)$ é obtido estritamente pela diferença de conjuntos:
\begin{equation}
    C_{viavel}(v) = \{1, 2, \dots, 9\} \setminus C_{indisponivel}(v)
\end{equation}

\subsection{Código Python: Grafo Explícito e Backtracking}
O trecho em Python a seguir implementa a classe \texttt{GrafoSudoku} e a busca por \textit{backtracking}:

\begin{lstlisting}[language=Python, caption={Algoritmo de Backtracking sobre Grafo Explícito em Python}]
class GrafoSudoku:
    def __init__(self):
        # 81 vertices indexados de 0 a 80
        self.vertices = list(range(81))
        # Lista de Adjacencias: N(v) para cada vertice
        self.adj = {v: set() for v in self.vertices}
        self._construir_arestas()

    def _obter_posicao(self, v):
        # Mapeia indice escalar v para (linha, coluna)
        return v // 9, v % 9

    def _construir_arestas(self):
        for u in self.vertices:
            r1, c1 = self._obter_posicao(u)
            b1 = (r1 // 3) * 3 + (c1 // 3)
            for v in self.vertices:
                if u != v:
                    r2, c2 = self._obter_posicao(v)
                    b2 = (r2 // 3) * 3 + (c2 // 3)
                    # Aresta de conflito se compartilham linha, coluna ou bloco 3x3
                    if r1 == r2 or c1 == c2 or b1 == b2:
                        self.adj[u].add(v)

    def obter_cores_validas(self, v, coloracao):
        # Operacao de conjuntos: {1..9} - {c(u) | u em N(v)}
        cores_vizinhos = {coloracao[u] for u in self.adj[v] if coloracao[u] != 0}
        return set(range(1, 10)) - cores_vizinhos

    def resolver_backtracking_grafo(self, coloracao, vertice_idx=0):
        # Condicao de Parada: Todos os 81 vertices coloridos
        if vertice_idx == 81:
            return True

        # Pula vertices com pistas fixas
        if coloracao[vertice_idx] != 0:
            return self.resolver_backtracking_grafo(coloracao, vertice_idx + 1)

        # Selecao de cores validas baseada estritamente em N(v)
        cores_possiveis = self.obter_cores_validas(vertice_idx, coloracao)
        for cor in sorted(cores_possiveis):
            coloracao[vertice_idx] = cor  # Atribui cor
            if self.resolver_backtracking_grafo(coloracao, vertice_idx + 1):
                return True
            coloracao[vertice_idx] = 0  # Backtrack (desfaz atribuicao)

        return False
\end{lstlisting}

\section{Algoritmo de Geração e Escavação com Solução Única}

Para gerar um enigma de Sudoku válido e matematicamente bem posto, é necessário assegurar a \textbf{unicidade da solução}. O processo se divide em duas etapas:
\begin{enumerate}
    \item Geração de uma $9$-coloração completa válida (tabuleiro resolvido aleatório).
    \item Processo de \textbf{escavação} (remoção de pistas) com teste de unicidade via poda computacional.
\end{enumerate}

\subsection{Código Python: Gerador de Enigmas de Sudoku}

\begin{lstlisting}[language=Python, caption={Gerador de Tabuleiros com Teste de Unicidade em Python}]
import random

class GeradorSudoku:
    def __init__(self):
        self.grafo = GrafoSudoku()

    def contar_solucoes(self, coloracao, vertice_idx=0, limite=2):
        if vertice_idx == 81:
            return 1

        if coloracao[vertice_idx] != 0:
            return self.contar_solucoes(coloracao, vertice_idx + 1, limite)

        total = 0
        cores_validas = self.grafo.obter_cores_validas(vertice_idx, coloracao)
        for cor in cores_validas:
            coloracao[vertice_idx] = cor
            total += self.contar_solucoes(coloracao, vertice_idx + 1, limite)
            coloracao[vertice_idx] = 0
            if total >= limite:  # Poda antecipada se encontrar mais de 1 solucao
                break
        return total

    def gerar_tabuleiro_completo(self):
        coloracao = [0] * 81
        def preencher_aleatorio(v_idx=0):
            if v_idx == 81:
                return True
            cores = list(self.grafo.obter_cores_validas(v_idx, coloracao))
            random.shuffle(cores)
            for cor in cores:
                coloracao[v_idx] = cor
                if preencher_aleatorio(v_idx + 1):
                    return True
                coloracao[v_idx] = 0
            return False

        preencher_aleatorio()
        return coloracao

    def criar_enigma(self, pistas_desejadas=30):
        tabuleiro = self.gerar_tabuleiro_completo()
        indices = list(range(81))
        random.shuffle(indices)

        removidos = 0
        alvo_remocao = 81 - pistas_desejadas

        for idx in indices:
            if removidos >= alvo_remocao:
                break
            cor_removida = tabuleiro[idx]
            tabuleiro[idx] = 0  # Escava a celula

            # Teste de Unicidade
            copia = list(tabuleiro)
            if self.contar_solucoes(copia) != 1:
                tabuleiro[idx] = cor_removida  # Reverte se houver mais de 1 solucao
            else:
                removidos += 1

        return tabuleiro
\end{lstlisting}

\section{Articulação Pedagógica: BNCC e Formação Docente}

\subsection{O Pensamento Computacional na BNCC}
A BNCC (Resolução CNE/CEB nº 1/2022) destaca o Pensamento Computacional como habilidade transversal indispensável na formação matemática contemporânea. A abordagem apresentada mobiliza diretamente os quatro pilares:
\begin{enumerate}
    \item \textbf{Decomposição:} Fracionamento do enigma global em restrições locais de vizinhança $N(v)$.
    \item \textbf{Abstração (EM13MAT501):} Modelagem do tabuleiro físico em um grafo não orientado $G=(V,E)$.
    \item \textbf{Reconhecimento de Padrões:} Mapeamento das regras de exclusão mútua como uma $k$-coloração própria.
    \item \textbf{Algoritmos (EM13MAT315 e EM13MAT405):} Implementação recursiva em linguagem Python com retrocesso investigativo.
\end{enumerate}

\subsection{Normatização ABNT}
A estrutura deste artigo atende às prescrições formais da \textbf{ABNT NBR 14724} para apresentação de trabalhos acadêmicos e \textbf{ABNT NBR 10520} para citações, garantindo o rigor metodológico indispensável à pesquisa em Educação Matemática.

\section{Considerações Finais}
A modelagem do Sudoku através da Teoria dos Grafos e Otimização Combinatória oferece uma ponte concreta entre o rigor matemático e a Ciência da Computação. O uso do algoritmo de \textit{backtracking} baseado em grafos explícitos em Python permite ao licenciando em Matemática vivenciar a transposição didática de conceitos abstratos, consolidando o Pensamento Computacional e fornecendo subsídios metodológicos para sua atuação docente.

\section*{Referências}
\begin{enumerate}
    \item ASSOCIAÇÃO BRASILEIRA DE NORMAS TÉCNICAS. \textbf{NBR 14724: Informação e documentação — Trabalhos acadêmicos — Apresentação}. Rio de Janeiro: ABNT, 2011.
    \item BRASIL. Ministério da Educação. \textbf{Base Nacional Comum Curricular (BNCC)}. Brasília: MEC, 2018.
    \item BRASIL. Conselho Nacional de Educação. \textbf{Resolução CNE/CEB nº 1/2022: Normas sobre Computação na Educação Básica}. Brasília: CNE, 2022.
    \item BONDY, J. A.; MURTY, U. S. R. \textbf{Graph Theory with Applications}. London: Macmillan, 1976.
    \item CORMEN, T. H. et al. \textbf{Algoritmos: Teoria e Prática}. 3. ed. Rio de Janeiro: Elsevier, 2012.
    \item WEST, D. B. \textbf{Introduction to Graph Theory}. 2. ed. Upper Saddle River: Prentice Hall, 2001.
\end{enumerate}

\end{document}
